\scanpage{305}
\mspage{\textcircled{4}}
la forme $A^*{}_*{}^{**}=(((A^*)_*)^*)^*$ \deleted{\unsure{etc, et de}} p.ex. -- dont les objets sont notés \deleted{par} par des lettres à étoile $K^{*}\unsure{{}^{*}_{*}}{}_*{}^{*}$ \ldots\ ]

Le rôle de $\Delta^*_{\deleted{\unsure{*}}}\in\mathrm{Ob}\,\widehat{\mathbb{\Delta}}^*\inserted{{}=\mathrm{Ob}\,\mathrm{Ss}^*}$ comme ``objet cosimplicial universel'' est explicité par la

\uline{Proposition} 1. Soit $A$ une catégorie où les $\underline{\lim}$ existent (resp $B$ une catégorie où les $\varprojlim$ existent). Alors
\[
 \begin{aligned}
 \text{a)}\quad A^*&\underset{\leftarrow}{\approx}\underline{\mathrm{Hom}}_{\substack{\text{comm}\\\text{à }\varinjlim}}(\deleted{\unsure{\widehat{\mathbb{\Delta}}}}\mathrm{Ss},A)\\
 \text{par}\quad F&\longmapsto F(\Delta^*)\\
 \text{b)}\quad B_*&\underset{\leftarrow}{\approx}\underline{\mathrm{Hom}}_{\substack{\text{comm}\\\text{à }\varprojlim}}(\deleted{\unsure{\widehat{\mathbb{\Delta}}}}\mathrm{Ss}^{o},\deleted{\unsure{A}}B)\\
 \text{par}\quad F&\longmapsto F((\Delta^*)^o)
 \end{aligned}
\]
(où $(\Delta^*)^o$ désigne l'objet simplicial de $\widehat{\mathbb{\Delta}}^o$ défini par l'objet cosimplicial $\Delta^*$ de $\widehat{\mathbb{\Delta}}$).

Bien sûr, b) n'est autre que a) appliqué à $A=B^o$, et a) signifie que le foncteur ``restriction à $\mathbb{\Delta}$''
\[
 \underline{\mathrm{Hom}}(\widehat{\mathbb{\Delta}},A)\longrightarrow\underline{\mathrm{Hom}}(\mathbb{\Delta},A)
\]
induit une équivalence sur la \unsure{ss-cat|s-cat.} \deleted{\unsure{des fonct|$\underline{\mathrm{Hom}}$}} \inserted{comm.\ à $\underline{\lim}$} du premier \unsure{mb.|membre} C'est vrai pour \uline{\unsure{tte|toute|les}} \deleted{(petite)} \inserted{\unsure{mt|toute}} catégorie $C$ -- pas seulement pour $\mathbb{\Delta}$, et bien connu. Le foncteur en sens inverse associe, à $K_*=\varphi:\mathbb{\Delta}\unsure{\longrightarrow|-} A$, le foncteur
\[
 X_*\longmapsto
 \underline{\lim}_{\Delta^n_*\in\mathbb{\Delta}/X_*}\varphi(\Delta^n_*)
 =\underline{\lim}_{\Delta^n_*\in\mathbb{\Delta}/X_*}K_{[n]}
\]
Dans le cas b) si on écrit
\[
 \underline{\mathrm{Hom}}(\widehat{\mathbb{\Delta}}^o,\deleted{\unsure{A}}B)\longrightarrow
 \underline{\mathrm{Hom}}(\mathbb{\Delta}^o,B)
\]
le foncteur en sens inverse associe à
\[
 K^*=\psi:\mathbb{\Delta}^o\unsure{\longrightarrow|-}B
\]
le foncteur contravariant
\[
 X^*\longmapsto\varprojlim_{\Delta^n_*\in(\mathbb{\Delta}/X_*)^o}\varphi(\Delta^n)
 =\varprojlim_{\Delta^n\in(\mathbb{\Delta}/X_*)^o}K^{[n]}
\]

\uline{Yoga} : nous nous intéresserons surtout \deleted{aux} (en première étape) aux invariants covariants attachés à $X_*$, à valeurs dans une cat. $A$, qui \deleted{\unsure{pour $X_*$ être|commutent aux $X_*$}} \inserted{commutent} aux $\underline{\lim}$ p.r. à $X_*$, -- ils sont \deleted{\unsure{donc}} exprimés (si $A$ est ``grande'' pour contenir les $\underline{\lim}$ \unsure{qqc|qqles|voulues}) par des éléments de $A^*$ i.e. des objets \uline{cosimpliciaux} de $A$. Pour les invariants contravariants en $X_*$, à val. dans une $B$, \deleted{\unsure{alors|dans}} nous attacherons en premier lieu à ceux qui \deleted{\unsure{con commutent}} aux $\varprojlim$ en $X_*$ ; pour $B$ stable par $\varprojlim$, ils correspondent \deleted{donc} aux objets de $B_*$, i.e. les objets $1/2$ simpliciaux de $B$. Si $B=(\mathrm{Ens})$, ils correspondent donc à des ens. $1/2$ simpliciaux, i.e. à des objets $K_*$ de $\deleted{\unsure{\widehat{\mathbb{\Delta}}}}\mathrm{Ss}$ lui-$=$.
% archive sheet number: 306
